% ── 5. Cloud Architecture Quality Attributes ("-ilities") ────────────────────
\section{Cloud Architecture Quality Attributes}

For this smart-home platform, the highest-priority quality attributes are security, availability, performance, scalability, fault tolerance, and durability. These matter most because the system must securely control household devices over the internet, remain responsive for actions such as unlocking doors or changing thermostat settings, and continue working as the number of homes and connected modules grows. The architecture supports these priorities by using authenticated device and user access, encrypted communication, managed cloud services, multi-Availability Zone data services, and loosely coupled components that can continue operating when individual failures occur. The main tradeoff is that this approach improves protection, responsiveness, and resilience, but it also increases reliance on AWS-managed services and leaves the initial single-region design more exposed to a full regional outage than a multi-region deployment would be.

Other important attributes include cost efficiency, maintainability, extensibility, interoperability, usability, and observability. These are important because the platform is intended for consumer households, must support modular device types over time, and is assumed to be operated by a relatively small team. The architecture supports them through usage-based managed services, centralized logging and monitoring, and a common cloud control layer that can accommodate new device types and protocol extensions without requiring a full redesign. The tradeoff is that these choices reduce operational overhead and make the system easier to evolve, but they also require careful cost monitoring, clear protocol design, and additional abstraction to keep the consumer experience simple.